% Source PDF pages 43--47; manual pages 2-14--2-18.
\subsection{Geodetic Coordinate System}\label{sec:geodetic}

The geodetic coordinate system is similar to the geocentric system, except that
position and velocity orientations are referenced to an ellipsoidal rather than a
spherical earth. The ellipsoidal model is more accurate, so geodetic coordinates are
used more often than geocentric coordinates. TAOS follows the method in
Reference~7 for the ellipsoidal geometry.

The earth's surface is modeled as an ellipsoid of revolution. In the equatorial plane,
defined by $\hat{x}_{\oplus}$ and $\hat{y}_{\oplus}$, the ellipsoid is a circle of
radius $R_{\oplus}$. The polar radius $R_p$ is less than the equatorial radius, so
points on the surface satisfy
\begin{equation}
  \frac{x_s^2}{R_{\oplus}^2}
  +\frac{y_s^2}{R_{\oplus}^2}
  +\frac{z_s^2}{R_p^2}=1.
  \label{eq:ellipsoid-surface}
\end{equation}

The ellipsoid is symmetric about the $\hat{z}_{\oplus}$ axis, so its geometry can be
analyzed in the $\hat{x}_{\oplus}$--$\hat{z}_{\oplus}$ plane, as shown in
\cref{fig:ellipsoidal-earth-geometry}.

\begin{figure}[htbp]
  \centering
  \resizebox{0.75\linewidth}{!}{\begin{tikzpicture}[scale=0.93,>=Latex,line cap=round,line join=round]
  \coordinate (O) at (0,0);
  \coordinate (Q) at (0,-1.32);
  \coordinate (S) at (3.02,1.25);
  \coordinate (P) at (4.12,2.20);
  \draw[line width=0.95pt] (O) ellipse [x radius=3.55,y radius=2.05];
  \draw[->,line width=0.9pt] (-4.15,0) -- (4.55,0) node[right] {$\hat{x}_{\oplus}$};
  \draw[->,line width=0.9pt] (0,-2.55) -- (0,2.75) node[above] {$\hat{z}_{\oplus}$};
  \draw[<->] (-3.98,0) -- (-3.98,2.05) node[midway,left] {$R_p$};
  \draw[<->] (-3.55,-2.38) -- (0,-2.38) node[midway,below] {$R_{\oplus}$};
  \draw[->,line width=1pt] (O) -- (P) node[midway,above left] {$\vec r$};
  \draw[line width=0.9pt] (Q) -- (P);
  \fill (Q) circle (1.4pt) node[below left] {$Q$};
  \fill (S) circle (1.4pt) node[above left] {$S$};
  \fill (P) circle (1.5pt) node[above right] {$P$};
  \draw[<->] ($(S)+(0.20,0.18)$) -- ($(P)+(0.20,0.18)$) node[midway,right] {$h$};
  \draw[<->] ($(Q)+(-0.22,0.15)$) -- ($(S)+(-0.22,0.15)$) node[midway,below right] {$N$};
  \draw[<->] ($(Q)+(0.18,0.08)$) -- ($(O)+(0.18,0.08)$) node[midway,right] {$Ne^2$};
  \draw[<->] (-0.48,0) -- (-0.48,-1.32) node[midway,left] {$z_i$};
  \draw[->] (1.20,0) arc[start angle=0,end angle=39,radius=1.20];
  \node at (1.38,0.50) {$\delta_{gd}$};
  \draw ($(S)+(-0.18,-0.18)$) -- ++(0.18,-0.18) -- ++(0.18,0.18);
\end{tikzpicture}
}
  \caption{Ellipsoidal earth geometry.}
  \label{fig:ellipsoidal-earth-geometry}
\end{figure}

The ellipse equation in this plane is
\begin{equation}
  R_{\oplus}^2=x_s^2+\frac{z_s^2}{1-e^2},
  \label{eq:ellipsoid-plane-equation}
\end{equation}
where $e$ is the eccentricity from \cref{eq:ellipsoid-parameter-relation}. Points on
the surface, such as $S$, satisfy this equation. A surface normal, such as the line
between $S$ and $Q$, has slope
\begin{equation}
  \tan\delta_{gd}=-\frac{d x_s}{d z_s}.
  \label{eq:geodetic-latitude-slope}
\end{equation}
Differentiating \cref{eq:ellipsoid-plane-equation} and substituting in
\cref{eq:geodetic-latitude-slope} gives
\begin{equation}
  \frac{z_s}{x_s}=(1-e^2)\tan\delta_{gd}.
  \label{eq:ellipsoid-surface-ratio}
\end{equation}
Solving \cref{eq:ellipsoid-plane-equation,eq:ellipsoid-surface-ratio} for $x_s$
produces
\begin{equation}
  x_s=\frac{R_{\oplus}\cos\delta_{gd}}
           {\sqrt{1-e^2\sin^2\delta_{gd}}}.
  \label{eq:ellipsoid-xs}
\end{equation}
From \cref{fig:ellipsoidal-earth-geometry}, point $S$ also satisfies
\begin{equation}
  x_s=N\cos\delta_{gd},
  \label{eq:ellipsoid-xs-normal}
\end{equation}
where $N$ is the distance from $S$ to the $\hat{z}_{\oplus}$ axis measured along the
surface normal. Consequently,
\begin{equation}
  N=\frac{R_{\oplus}}
          {\sqrt{1-e^2\sin^2\delta_{gd}}}.
  \label{eq:ellipsoid-normal-distance}
\end{equation}
Substitution into \cref{eq:ellipsoid-surface-ratio} gives
\begin{equation}
  z_s=(N-e^2N)\sin\delta_{gd}
     =N(1-e^2)\sin\delta_{gd}.
  \label{eq:ellipsoid-zs}
\end{equation}
Thus the distance between the $z$-axis intercept of the normal, point $Q$, and its
$x$-axis intercept is $Ne^2$.

\taossubhead{Geodetic Position}

The ellipsoidal relationships locate the vehicle relative to the earth's surface. In
\cref{fig:ellipsoidal-earth-geometry,fig:geodetic-coordinate-geometry}, the vehicle
is at point $P$, given by $\vec{r}$, and point $S$ is on the surface directly below it.

\noindent\begin{tabularx}{\textwidth}{@{}>{\ttfamily}l>{$\displaystyle}l<{$}X@{}}
  alt & h &
    Altitude: distance from the vehicle's center of mass to the ellipsoid, measured
    along the surface normal from $S$ to $P$. \\
  long & \lambda &
    Longitude: the angle between the ECFC $x$ axis and the projection of
    $\vec{r}$ onto the equatorial plane, positive eastward. \\
  latgd & \delta_{gd} &
    Geodetic latitude: the angle between the local geodetic north/down geometry
    and the equatorial plane; equivalently, the angle that the ellipsoid surface
    normal makes with the equatorial plane. It is positive northward. \\
\end{tabularx}

\begin{figure}[htbp]
  \centering
  \resizebox{0.72\linewidth}{!}{\begin{tikzpicture}[scale=0.90,>=Latex,line cap=round,line join=round]
  \coordinate (O) at (0,0);
  \coordinate (Q) at (0,-1.38);
  \coordinate (S) at (2.25,1.05);
  \coordinate (P) at (3.15,2.10);
  \coordinate (E) at (3.15,0);
  \draw[line width=0.95pt] (O) ellipse [x radius=3.10,y radius=2.05];
  \draw[->,line width=0.9pt] (O) -- (-2.65,-1.55) node[below left] {$\hat{x}_{\oplus}$};
  \draw[->,line width=0.9pt] (O) -- (4.20,0) node[right] {$\hat{y}_{\oplus}$};
  \draw[->,line width=0.9pt] (O) -- (0,2.85) node[above] {$\hat{z}_{\oplus}$};
  \draw[->,line width=1pt] (O) -- (P) node[midway,above left] {$\vec r$};
  \draw[line width=0.9pt] (Q) -- (P);
  \fill (Q) circle (1.3pt) node[below] {$Q$};
  \fill (S) circle (1.3pt) node[above left] {$S$};
  \fill (P) circle (1.5pt) node[above right] {$P$};
  \draw[dashed] (P) -- (E);
  \draw[line width=0.75pt] (O) -- (E);
  \draw[<->] (3.52,0) -- (3.52,2.10) node[midway,right] {$\vec r\cdot\hat z_{\oplus}$};
  \draw[<->] (0,-1.70) -- (0,-0.05) node[midway,left] {$z_i$};
  \node[align=left,right] at (3.35,-0.85) {$\tilde{x}=$ projection of $\vec r$ onto\\the equatorial plane};
  \draw[->] (0.70,-0.23) arc[start angle=-18,end angle=-142,radius=0.72];
  \node at (-0.05,-0.70) {$\lambda$};
  \draw[->] (1.25,0) arc[start angle=0,end angle=35,radius=1.25];
  \node at (1.50,0.52) {$\delta_{gd}$};
  \draw ($(S)+(-0.12,-0.15)$) -- ++(0.17,-0.16) -- ++(0.17,0.17);
\end{tikzpicture}
}
  \caption{Geodetic coordinate-system geometry.}
  \label{fig:geodetic-coordinate-geometry}
\end{figure}

Longitude is identical in the geocentric and geodetic systems and is computed from
\cref{eq:longitude-from-ecfc}. The remaining ECFC/geodetic position relationships
are
\begin{equation}
  \vec{r}\cdot\hat{x}_{\oplus}=\tilde{x}\cos\lambda,
  \label{eq:geodetic-position-x}
\end{equation}
\begin{equation}
  \vec{r}\cdot\hat{y}_{\oplus}=\tilde{x}\sin\lambda,
  \label{eq:geodetic-position-y}
\end{equation}
\begin{equation}
  \vec{r}\cdot\hat{z}_{\oplus}
    =\bigl[N(1-e^2)+h\bigr]\sin\delta_{gd},
  \label{eq:geodetic-position-z}
\end{equation}
where
\begin{equation}
  \tilde{x}=(N+h)\cos\delta_{gd}.
  \label{eq:geodetic-equatorial-projection}
\end{equation}
The function \texttt{p\_geodetic\_to\_ecfc} evaluates these expressions when the
geodetic quantities are known. In the reverse direction, altitude and geodetic
latitude cannot be isolated directly because $\delta_{gd}$ appears
transcendentally. An iterative solution is therefore required.

Extending $\hat{z}_{gd}$ to the ECFC $z$ axis gives the intercept distance
\begin{equation}
  z_i=Ne^2\sin\delta_{gd}
     =\frac{R_{\oplus}e^2\sin\delta_{gd}}
            {\sqrt{1-e^2\sin^2\delta_{gd}}}.
  \label{eq:geodetic-axis-intercept}
\end{equation}
For an initial estimate, the earth is treated as nearly spherical so that
\begin{equation}
  \sin\delta_{gd}\approx\sin\delta_{gc}
    =\frac{\vec{r}\cdot\hat{z}_{\oplus}}{\lVert\vec{r}\rVert}.
  \label{eq:initial-geodetic-latitude-estimate}
\end{equation}
This approximation is poor when $\lVert\vec{r}\rVert\ll R_{\oplus}$, but it is
sufficient to start the iteration. The initial intercept estimate is
\begin{equation}
  z_i=R_{\oplus}e^2\sin\delta_{gc}.
  \label{eq:initial-axis-intercept}
\end{equation}
The following sequence is then repeated:
\begin{equation}
  z_d=\vec{r}\cdot\hat{z}_{\oplus}+z_i,
  \label{eq:geodetic-iteration-zd}
\end{equation}
\begin{equation}
  N+h=\sqrt{\tilde{x}^{2}+z_d^{2}},
  \label{eq:geodetic-iteration-altitude}
\end{equation}
\begin{equation}
  \sin\delta_{gd}=\frac{z_d}{N+h},
  \label{eq:geodetic-iteration-latitude}
\end{equation}
\begin{equation}
  N=\frac{R_{\oplus}}
          {\sqrt{1-e^2\sin^2\delta_{gd}}},
  \label{eq:geodetic-iteration-normal}
\end{equation}
\begin{equation}
  z_i=Ne^2\sin\delta_{gd}.
  \label{eq:geodetic-iteration-intercept}
\end{equation}
Iteration terminates when successive $z_i$ values differ by less than $10^{-8}$.
Typical applications achieve approximately \SI{0.01}{\foot} accuracy in fewer than
three iterations; no more than 25 iterations are performed. Altitude is obtained
from \cref{eq:geodetic-iteration-altitude}, and geodetic latitude from
\cref{eq:geodetic-iteration-latitude}.

\taossubhead{Geodetic Velocity}

\noindent\begin{tabularx}{\textwidth}{@{}>{\ttfamily}l>{$\displaystyle}l<{$}X@{}}
  vel & \lVert\vec{V}_{\oplus}\rVert &
    Magnitude of the earth-relative velocity vector. \\
  gamgd & \gamma_{gd} &
    Geodetic vertical flight-path angle: the angle between the velocity vector and
    the local geodetic horizon plane, positive away from the earth's surface. \\
  psigd & \psi_{gd} &
    Geodetic horizontal flight-path, or heading, angle: the angle between local
    north and the projection of velocity onto the local geodetic horizon plane,
    positive clockwise from north so that $90^{\circ}$ is east. \\
\end{tabularx}

The earth-relative speed is the same in geocentric and geodetic coordinates. The
geodetic flight-path angles are defined by
\begin{equation}
  \vec{V}_{\oplus}\cdot\hat{x}_{gd}
    =\lVert\vec{V}_{\oplus}\rVert
      \cos\gamma_{gd}\cos\psi_{gd},
  \label{eq:geodetic-velocity-x}
\end{equation}
\begin{equation}
  \vec{V}_{\oplus}\cdot\hat{y}_{gd}
    =\lVert\vec{V}_{\oplus}\rVert
      \cos\gamma_{gd}\sin\psi_{gd},
  \label{eq:geodetic-velocity-y}
\end{equation}
\begin{equation}
  \vec{V}_{\oplus}\cdot\hat{z}_{gd}
    =-\lVert\vec{V}_{\oplus}\rVert\sin\gamma_{gd}.
  \label{eq:geodetic-velocity-z}
\end{equation}
Therefore,
\begin{equation}
  \sin\gamma_{gd}
    =-\frac{\vec{V}_{\oplus}\cdot\hat{z}_{gd}}
            {\lVert\vec{V}_{\oplus}\rVert},
  \label{eq:geodetic-gamma-from-ecfc}
\end{equation}
\begin{equation}
  \tan\psi_{gd}
    =\frac{\vec{V}_{\oplus}\cdot\hat{y}_{gd}}
           {\vec{V}_{\oplus}\cdot\hat{x}_{gd}}.
  \label{eq:geodetic-psi-from-ecfc}
\end{equation}
The function \texttt{v\_ecfc\_to\_geodetic} evaluates these equations using the
unit vectors in
\cref{eq:local-geodetic-x-unit-vector,eq:local-geodetic-y-unit-vector,eq:local-geodetic-z-unit-vector}.
They are identical in form to the geocentric equations, except that the local geodetic
horizon replaces the local geocentric horizon.

The reverse calculation, performed by \texttt{v\_geodetic\_to\_ecfc}, is
\begin{equation}
\begin{aligned}
  \vec{V}_{\oplus}\cdot\hat{x}_{\oplus}
    ={}&-(\vec{V}_{\oplus}\cdot\hat{x}_{gd})
         \cos\lambda\sin\delta_{gd}
       -(\vec{V}_{\oplus}\cdot\hat{y}_{gd})\sin\lambda\\
      &-(\vec{V}_{\oplus}\cdot\hat{z}_{gd})
         \cos\lambda\cos\delta_{gd},
\end{aligned}
\label{eq:geodetic-to-ecfc-x}
\end{equation}
\begin{equation}
\begin{aligned}
  \vec{V}_{\oplus}\cdot\hat{y}_{\oplus}
    ={}&-(\vec{V}_{\oplus}\cdot\hat{x}_{gd})
         \sin\lambda\sin\delta_{gd}
       +(\vec{V}_{\oplus}\cdot\hat{y}_{gd})\cos\lambda\\
      &-(\vec{V}_{\oplus}\cdot\hat{z}_{gd})
         \sin\lambda\cos\delta_{gd},
\end{aligned}
\label{eq:geodetic-to-ecfc-y}
\end{equation}
\begin{equation}
  \vec{V}_{\oplus}\cdot\hat{z}_{\oplus}
    =(\vec{V}_{\oplus}\cdot\hat{x}_{gd})\cos\delta_{gd}
     -(\vec{V}_{\oplus}\cdot\hat{z}_{gd})\sin\delta_{gd}.
  \label{eq:geodetic-to-ecfc-z}
\end{equation}
These are the same transformations as
\cref{eq:geocentric-to-ecfc-x,eq:geocentric-to-ecfc-y,eq:geocentric-to-ecfc-z},
with the geodetic subscript replacing the geocentric subscript.
